\section{Introduction}
\label{sec:intro}

Congressional redistricting requires partitioning a weighted planar graph
into $k$ connected, population-balanced parts while minimising total edge
cut (shared boundary perimeter).
The algorithmic literature has long recognised that this problem is NP-hard
in general~\citep{garey1979computers,garey1976simplified}, motivating the
use of heuristic partitioners such as METIS~\citep{karypis1998}.
However, the precise complexity status of the \emph{redistricting variant}
--- with its planarity constraint, population-balance requirement, and
boundary-length edge weights --- has not been rigorously established.

Moreover, the empirical scaling behaviour of METIS recursive bisection on
real census-tract graphs across three decades of census data has not been
characterised.  For the Districting Integrity Act (DIA) to be legally
defensible, practitioners and courts need to know: (1) is the optimisation
problem NP-hard even for the planar case? (2) what approximation guarantee
does METIS offer? (3) does the algorithm scale to large states and future
Census enumerations?

\subsection*{Contributions}

This paper makes four contributions:

\begin{enumerate}
  \item We establish NP-hardness of the balanced $k$-cut redistricting
    problem on planar graphs by reduction from Connected Planar Graph
    Bisection~\citep{dyer1985partitioning}, correctly accounting for the
    contiguity (connectivity) requirement of the redistricting problem.
  \item We clarify the approximation-theoretic landscape: no formal
    approximation ratio is known for METIS's heuristic; the best
    polynomial-time bound ($O(\sqrt{\log n})$ via ARV~\citep{arora2009expander})
    applies to Sparsest Cut, not to balanced connected $k$-partition.
    We characterise METIS's empirical quality instead.
  \item We provide a runtime analysis: $O(\texttt{niter} \cdot n \log k)$
    time with $O(n \log k)$ space ($\texttt{niter} = 100$ in experiments).
  \item We fit the empirical scaling law $T = a \cdot n^b$ to $k > 1$
    states across three census years, finding $b = 1.07 \pm 0.03$
    (near-linear); single-district states are excluded from the fit.
\end{enumerate}

\subsection*{Why complexity matters for redistricting law}

The DIA's algorithm specification needs a complexity-theoretic foundation
for two reasons.  First, if the redistricting problem were polynomial-time
solvable, the DIA would be obligated to specify an exact algorithm rather
than a heuristic.  Our NP-hardness result shows that no such exact
algorithm can exist unless P = NP (a hypothesis universally believed to
be false, making the result effectively unconditional in practice).
Second, while no formal approximation guarantee is established for METIS's
heuristic, the empirical characterisation in Section~\ref{sec:runtime}
shows that METIS consistently finds solutions within 3\% of the best
available partition~\citep{b7paper}, providing strong practical evidence
for legal challenges alleging systematic suboptimality.
Importantly, the hardness result concerns the \emph{optimisation} problem;
a valid redistricting plan always exists and can be found efficiently ---
only finding the \emph{optimal} one is intractable.

\subsection*{Paper structure}

Section~\ref{sec:formulation} formally defines the redistricting problem,
establishes NP-hardness, and characterises the approximation landscape.
Section~\ref{sec:runtime} analyses runtime complexity, empirical scaling,
and space requirements.
Section~\ref{sec:discussion} discusses implications for the DIA.
Section~\ref{sec:conclusion} concludes.
